% ============================================================
% COMPOSER'S NOTE § 4 — Making the Score
% ============================================================
% The organizing metaphor: making an instrument — not
% vessel/voyage, but the act of construction itself.
% Sound is both tool AND material, instrument AND subject.
% Two temporalities named explicitly (polytonality):
%   — Geological/elemental: earth → water → body
%   — Historical/linear: past → speculative future → present
% Both true simultaneously.
% Brief sketch of ST framework: URET cycle + four domains,
% one or two sentences each. Just enough to have ears.
% Julius Eastman's Femenine named as structural model:
%   drone (perceptual constant) + developmental cell (ST).
%
% Key texts:
%   Eastman — Femenine (structural model)
%   Manning & Massumi — Thought in the Act
%   Loveless — How to Make Art at the End of the World
% ============================================================

\section*{Making the Score: On the Structure of This Dissertation}

The tools through which we create and perceive sound carry within them specific cultural orientations. Digital audio systems, while presented as universal, encode the temperaments and assumptions of their origins. Equal temperament tuning, grid-based sequencing, and Eurocentric standards of fidelity organize perception according to a narrow acoustic worldview. These defaults do more than shape what can be heard; they shape what can be imagined.

When Ruha Benjamin warns that ``the road to inequity is paved with technical fixes'' \parencite[p.~13]{Benjamin_2019a}, she identifies why surface-level interventions often fail to address the foundational issue: technical fixes assume problems lie in biased implementations of otherwise neutral procedures---maintaining what Sylvia Wynter calls the over-representation of ``Man,'' the Western bourgeois conception that has ``positioned itself as equivalent to humanity itself'' through supposedly universal logics \parencite{Wynter_2003}. \textit{Speculocultural Technopoiesis} responds to this condition by proposing a form of sonic research-creation grounded in modification. Through iterative cycles of making, listening, and reflection, it explores how culturally informed approaches to technology can generate new sonic and epistemic possibilities.

In this view, modification functions as a form of \textit{transduction}: speculative cultural ideas are converted into material form through the redesign of software, instruments, and interfaces. The resulting shifts in perception---what this framework calls \textit{refraction}---make audible relationships that were previously muted by the defaults of dominant design. Black creative practices have consistently demonstrated how the same tools produce fundamentally different possibilities when approached through diametric ways of knowing. These practices affirm that tools cannot be neutral containers. They function as systems that morph the shape and function of their mediated perception when operated through distinct hands.

The framework of \textit{Speculocultural Technopoiesis} (ST) proceeds through four movements that repeat iteratively across each work:

\begin{itemize}
    \item \textbf{Unsettling} --- surfacing the friction between cultural knowledge and the tool's embedded assumptions; identifying what must change.
    \item \textbf{Remixing} --- materially modifying tools, interfaces, and algorithms to encode different cultural logics.
    \item \textbf{Embodying} --- testing cultural authenticity and experiential fit through practice; validating what the modification enables.
    \item \textbf{Transmitting} --- documenting, naming, and sharing methods and code so that knowledge circulates with its cultural reasoning intact.
\end{itemize}

Four cross-cutting domains---\textit{Fugitive Speculation}, \textit{Material Memory}, \textit{Sensory Worlding}, and \textit{Intimate Witnessing}---guide design decisions and evaluative attention in every iteration.

This score is organized around two simultaneous temporalities. One is elemental and geological: earth (Ground) gives way to water (Hold) gives way to the body (Touch). The other is historical and speculative: the works move from the buried past of the Buffalo Soldier to the suspended middle passage of cryosleep to the present-tense politics of Black hair. Both are always true at once, the way a chord is never just one note.

\adnote{The Julius Eastman / \textit{Femenine} reference (drone as perceptual constant + ST as developmental cell) is flagged in the original scaffolding---this structural analogy is strong and should be developed here. Also, Manning \& Massumi and Loveless are cited in the scaffolding but not yet brought in. This section needs expansion on those fronts when you're ready.}
